% !TeX root = ../../Book3.tex
\section{正规族与省略值}
\label{b3:sec:1403}

省略值将像限制在一个带穿孔的球面。正规性可以通过球面导数刻画；若正规性失败，重标度会制造一条整个平面上的非恒常极限曲线，与穿孔球面的圆盘覆盖矛盾。

\subsection{省略值}
\label{b3:subsec:140301}

亚纯函数省略一个球面值，指它从不取到该值，极点被解释为取无穷远。省略三个不同球面值的族可用同一个 Möbius 变换把这三个值化为0、1、无穷远。全纯函数本已省略无穷远，所以再省略两个不同有限值便属于此情形。

\ProseParagraphBreak
球面导数定义为 \(f^\#=|f'|/(1+|f|^2)\)，极点处用 \(1/f\) 同式定义，两个表达一致。它测量到球面长度元 \(|\mathrm{d}w|/(1+|w|^2)\) 的导数大小，不是普通导数。\(f_n=n\) 的普通值无界，但球面导数0且球面趋于无穷远。

\subsection{正规性}
\label{b3:subsec:140302}

\term[marty-panbie]{Marty 判别}[Marty's criterion]以球面导数表达正规性。

\begin{theorem}[Marty 判别]\label{b3:thm:14-marty}
亚纯函数族球面正规，当且仅当其球面导数在每个紧集有共同上界。
\end{theorem}
\begin{proof}
若导数局部共同有界，球面两点的弦距离不超过连接像路径的球面长度，故在小凸盘上得到共同 Lipschitz 界。球面紧，对可数稠密点对角抽取，再用等度连续性与有限网格得到紧集一致子列。连续极限在每个点可选一个避开无穷远的球面坐标；充分晚成员也在此坐标有界全纯，Weierstrass 定理使极限全纯。所以极限为亚纯或球面常值，得到正规性。

反向若球面导数在紧集无界，取 \(f_n\)、\(z_n\rightarrow a\) 使 \(f_n^\#(z_n)\rightarrow\infty\)。正规性给球面一致子列。在 \(a\) 的极限值附近选一个球面旋转坐标，把该值送到有限点；在小盘中成员普通一致收敛且共同有界，Cauchy 导数估计给共同导数界。球面旋转保持球面导数，产生矛盾。极限恒为无穷远时同样选倒数坐标。
\end{proof}

\ProseParagraphBreak
Marty 判别把欧氏导数改成球面导数，能在极点处继续使用。若 \(g=1/f\)，则 \(|g'|/(1+|g|^2)=|f'|/(1+|f|^2)\)，两张球面坐标给相同数值。以 \(f_n(z)=nz\) 为例，0处球面导数为 \(n\)，族在那里不正规；但在任一避开0的紧集上，\(f_n\) 球面趋于无穷，正规性完全可以恢复。正规是局部性质，不能因为某一点失败就忽略其他子域的正常极限。

\ProseParagraphBreak
另一方面，\(f_n(z)=n\) 是处处球面导数为0的正规族，虽没有有限的 Euclid 局部有界性。第四章有限全纯极限的定义与此处球面正规定义使用不同紧目标，必须分清。把极限 \(\infty\) 加入球面以后，常数逃逸成为紧致现象；却没有因此把任何无界族都变成正规，\(nz\) 已给反例。

\subsection{紧性}
\label{b3:subsec:140303}

\term[zalcman-chongsuofang-yinli]{Zalcman 重缩放引理}[Zalcman's rescaling lemma]把不正规性转成非常值整平面极限。

\begin{lemma}[Zalcman 重缩放引理]\label{b3:lem:14-zalcman-rescaling}
若亚纯族在 \(a\) 附近不球面正规，则存在成员 \(f_n\)、\(a_n\rightarrow a\)、\(r_n\downarrow0\)，使 \(f_n(a_n+r_n\zeta)\) 在复平面紧集上球面趋于一个非恒常亚纯函数。
\end{lemma}
\begin{proof}
由 Marty，选择缩小圆盘半径 \(s_n\downarrow0\) 及成员，使 \(|z-a|<s_n/2\) 内某点的球面导数超过 \(n^2/s_n\)。在闭盘上最大化 \((s_n-|z-a|)f_n^\#(z)\)，取最大点 \(a_n\)，记 \(d_n=s_n-|a_n-a|\)、\(r_n=1/f_n^\#(a_n)\)。则 \(d_n/r_n\rightarrow\infty\)，且 \(r_n\rightarrow0\)。当 \(|\zeta|\le d_n/(2r_n)\)，最大性给重标度函数球面导数不超过2，而在0恰为1。对越来越大的圆盘用 Marty 与对角抽取，得到全平面亚纯极限。局部有限球面坐标中的 Cauchy 导数收敛使极限在0球面导数为1，故不恒常。
\end{proof}
重标度既选中心又选尺度；只令 \(f_n(r_nz)\) 而不寻找导数最大处，极限可能常值，无法反映正规性失败。

\ProseParagraphBreak
重缩放的目的并非单纯放大图形，而是同时固定一个非零导数和越来越大的定义域。\(f_n(z)=nz\) 在0处取尺度 \(1/n\)，得到 \(g_n(\zeta)=\zeta\)，定义域半径扩大 \(n\) 倍，极限为非常值整函数。\(f_n(z)=e^{nz}\) 在0处球面导数为 \(n/2\)，取尺度 \(2/n\) 得 \(g_n=e^{2\zeta}\)，极限仍非常值，省略0与无穷远而不省略第三个球面值。这个例子也显示三省略值定理中的数目尖锐。

\ProseParagraphBreak
如果仅固定原点而族的坏行为发生在移动点，那里导数可能不反映最大尺度。Zalcman 选择点距边界乘球面导数的最大值，保证新中心不会太靠近旧边界，并使缩放后的导数在每个固定圆盘受统一界控制。三个任务分别是捕获坏行为、扩大定义域和取得紧极限，缺少任何一个都不能完成矛盾论证。

\subsection{广义 Montel 定理}
\label{b3:subsec:140304}

\begin{theorem}[省略值 Montel 定理]\label{b3:thm:14-montel-omitted}
省略同三个不同球面值的亚纯函数族球面正规。省略同两个不同有限值的全纯函数族也球面正规，其极限为全纯函数或恒为无穷远。
\end{theorem}
\begin{proof}
若不正规，重标度给非恒常全平面亚纯极限 \(g\)。对每个省略值选局部球面坐标，Hurwitz 无零定理说明极限要么仍省略该值，要么恒取该值；非恒常排除后者。因此 \(g:\mathbb C\rightarrow\widehat{\mathbb C\kern0pt}\setminus\{a,b,c\}\)。三穿孔球面由第13章覆盖输入有圆盘万有覆盖，单连通平面使 \(g\) 提升为全纯 \(G:\mathbb C\rightarrow\mathbb D\)。Liouville 使 \(G\) 恒常，进而 \(g\) 恒常，矛盾。全纯列的球面极限若非恒为无穷远，不会产生内点极点，已由第4章 Hurwitz 倒数证明。
\end{proof}
这里补全第4章预告的省略值版本，没有在统一化或 Riemann 极值证明中回用它，因而依赖路线无循环。单个有限省略值不足，\(e^{nz}\) 在含0的盘给反例。

\ProseParagraphBreak
省略值可以随函数变化，但须避免相互碰撞。若每个 \(f_n\) 省略三个球面点 \(a_n,b_n,c_n\)，且它们两两球面距离有统一正下界，则先取子列使三点收敛到三个不同点。将它们送到 \(0,1,\infty\) 的 Möbius 变换及其逆随参数连续，在球面上一致收敛到非退化变换；正规三省略值定理应用于变换后的族，再复合逆变换即可。此论证也给族的任意序列子列判别。

\ProseParagraphBreak
若允许碰撞，结论失败。例如 \(f_n(z)=nz\) 在单位盘省略 \(2n,3n,4n\)，三个点球面上都趋向无穷，族仍在0不正规。“每个函数都有三个省略值”若不记录这三个值共同固定或有统一分离，不能当作 Montel 定理的假设。
